% Originally: content/week-08.tex, \subsection{Change of variables} --- promoted to a section

\section{Change of variables}
\label{sec:change_of_variables}

\begin{theorem}[Change of variables]
\label{thm:change_of_variables}
Suppose $U,V \subseteq \mathbb{R}^n$ are open sets and $\Phi : U \to V$ is a
$C^1$-diffeomorphism. If $A \subseteq U$ is Jordan measurable with $\overline{A} \subseteq U$,
and $f : A \to \mathbb{R}$ is continuous, then
\[\int_A f(x)\,dx
  = \int_{\Phi(A)} \frac{f\circ\Phi^{-1}(y)}{\big|\det \Jac\Phi\big(\Phi^{-1}(y)\big)\big|}\,dy.\]
\end{theorem}

\begin{ainote}
Corsin uses \newterm{Jordan measurable} (\germanterm{Jordan-messbar}) here without defining it;
the definition is in the lecture (ch.~13). For this chapter it is enough to know that boxes,
and the bounded regions cut out by finitely many smooth curves that appear in the examples, are
all Jordan measurable.

On the shape of the formula: it is often quoted the other way round, as
$\int_{\Phi(A)} f = \int_A (f\circ\Phi)\,\lvert\det\Jac\Phi\rvert$. The version above is the same
statement with $\Phi$ pushed to the other side, which is why the determinant appears
\emph{divided} rather than multiplied. Which form is convenient depends on whether it is $\Phi$
or $\Phi^{-1}$ that you can write down --- in \cref{ex:change_of_variables_domain} below it is
$\Phi$, so the dividing form is the one that saves work.
\end{ainote}

% Supplement: Sascha Brack/Class Notes/Week_08_Notes_Friday.pdf, pp. 2--3
\begin{example}[Change of variables on a diamond]
\label{ex:change_of_variables_diamond}
Compute $I := \int_E e^{x+y}(x-y) \mathop{d\text{vol}}(x,y)$, where $E$ is the region bounded by the lines $y=x$, $y=-x+2$, $y=x-2$, and $y=-x$.

Writing out the inequalities, $E = \{ y \leq x \} \cap \{ y \leq -x+2 \} \cap \{ y \geq -x \} \cap \{ y \geq x-2 \}$.
Rearranging these gives $0 \leq x-y \leq 2$ and $0 \leq x+y \leq 2$.
This suggests the substitution $u := x+y$ and $v := x-y$. The new domain is the square $V = [0,2] \times [0,2]$.
The inverse map is $\Phi^{-1}(u,v) = (x,y) = \big(\frac{u+v}{2}, \frac{u-v}{2}\big)$.
Its Jacobian matrix is $\Jac\Phi^{-1}(u,v) = \begin{pmatrix} 1/2 & 1/2 \\ 1/2 & -1/2 \end{pmatrix}$, so $\big|\det\Jac\Phi^{-1}\big| = \big|-\frac{1}{4} - \frac{1}{4}\big| = \frac{1}{2}$.
By the change of variables formula:
\[I = \int_0^2 \int_0^2 e^u v \left|\det\Jac\Phi^{-1}\right| \,du\,dv = \frac{1}{2} \int_0^2 \int_0^2 e^u v \,du\,dv = \frac{1}{2} \left[ e^u \right]_0^2 \left[ \frac{v^2}{2} \right]_0^2 = \frac{1}{2} (e^2-1) \cdot 2 = e^2-1.\]
\end{example}

\begin{example}[Improper integral via polar coordinates]
\label{ex:improper_integral_polar}
Compute $\int_{\mathbb{R}^2} \frac{1}{(1+x^2+y^2)^2} \mathop{d\text{vol}}(x,y)$.
This integral is improper because the domain is unbounded. Since the integrand is non-negative and continuous, its Riemann improper integral is well-defined, and we can compute it using polar coordinates $x = r\cos\theta, y = r\sin\theta$.
The Jacobian determinant is $r$. The domain $\mathbb{R}^2$ corresponds to $r \in [0, \infty)$ and $\theta \in [0, 2\pi)$.
\[
\int_{\mathbb{R}^2} \frac{1}{(1+x^2+y^2)^2} \,dx\,dy
= \int_0^{2\pi} \int_0^\infty \frac{r}{(1+r^2)^2} \,dr\,d\theta.
\]
Substituting $u = r^2$, $du = 2r\,dr$, the inner integral is $\frac{1}{2} \int_0^\infty \frac{1}{(1+u)^2} \,du = \frac{1}{2} \left[ \frac{-1}{1+u} \right]_0^\infty = \frac{1}{2}$.
The outer integral then gives $\frac{1}{2} \int_0^{2\pi} d\theta = \pi$.
\end{example}

% Supplement: Sascha Brack/Class Notes/Week_08_Notes_Monday.pdf, pp. 8--9
\begin{example}[Volume of a sphere]
\label{ex:volume_sphere}
Let $A := \{(x,y,z) \in \mathbb{R}^3 : x^2+y^2+z^2 \leq 1\}$ be the unit ball. Given that $A$ is Jordan measurable, its volume is $\mu(A) = \int_A 1 \mathop{d\text{vol}}(x,y,z)$.
We can compute this using spherical coordinates:
\[\Phi(r,\theta,\varphi) = \begin{pmatrix} r\sin\theta\cos\varphi \\ r\sin\theta\sin\varphi \\ r\cos\theta \end{pmatrix}.\]
Up to a Jordan null set (which does not change the volume), the domain $A$ corresponds to the box $V = (0,1) \times (0,\pi) \times (-\pi,\pi)$.
The Jacobian determinant is $\det\Jac\Phi(r,\theta,\varphi) = r^2\sin\theta$, which is positive on this domain.
Applying the change of variables formula and Fubini's theorem (which allows slicing up the integral):
\begin{align*}
\mu(A) &= \int_{\Phi^{-1}(A)} \big|\det\Jac\Phi(r,\theta,\varphi)\big| \,dr\,d\theta\,d\varphi \\
&= \int_{-\pi}^\pi \int_0^\pi \int_0^1 r^2\sin\theta \,dr\,d\theta\,d\varphi \\
&= \left(\int_{-\pi}^\pi d\varphi\right) \left(\int_0^\pi \sin\theta \,d\theta\right) \left(\int_0^1 r^2 \,dr\right) \\
&= (2\pi) \cdot (2) \cdot \left(\frac{1}{3}\right) = \frac{4\pi}{3}.
\end{align*}
\end{example}

% Source: Corsin Nick/Class Notes/Week 8.pdf, p. 7


% --- exercises relocated here from the dissolved sheet 8 ---

% Originally: content/week-08.tex, \exercisesheet{8}, problem 8.5
% Source: exercises/Ex8_Analysis2_eng.pdf

% Source: exercises/Ex8_Analysis2_eng.pdf, pp. 1--2
\begin{exercise}[8.5 --- Change of variables and Jacobians \textnormal{(important)}]
\label{ex:8.5}
For each of the following domains and change of variables find the jacobian and the appropriate
transformed domain. There is no need to actually compute the integrals!
\begin{enumerate}[label=\textbf{\arabic*.}]
  \item $A := \{(x_1,x_2) \in \mathbb{R}^2 : x_1 > 0,\ x_2 < x_1,\ 1 < x_1^2+x_2^2 < 4\}$ and
    $x_1 = r\cos\theta$, $x_2 = r\sin\theta$. Using the change of variables formula, complete
    the dots in the following formula
    \[\int_A x_1^2\sin(x_2)\,dx_1dx_2 = \int_{\dots}\dots\,dr\,d\theta\]
  \item $B := \{(x,y) \in \mathbb{R}^2 \mid 1 < xy < 2,\ x^2 < y < 2x^2\}$ and $u := xy$,
    $v := x^2$. Using the change of variables formula, complete the dots in the following
    formula
    \[\int_B y^2e^{-xy}\,dx\,dy = \int_{\dots}\dots\,du\,dv\]
  \item $C := \{(x,y,z) \in \mathbb{R}^3 \mid 1 < z-2y < 2,\ 0 < z < 1,\ -2 < 3x+y+z < 0\}$ and
    $u := z$, $v := z-2y$, $w := 3x+y+z$. Using the change of variables formula, complete the
    dots in the following formula
    \[\int_C xyz\,dx\,dy = \int_{\dots}\dots\,du\,dv\,dw.\]
\end{enumerate}

% Supplement: Sascha Brack/Ex Sheet Hints/Ex8_Analysis2_hints.pdf, p. 1
\textit{Sascha Brack's hint} on \textbf{8.5.2:} first pin down the signs. The constraints force
$x \neq 0$ and $y \neq 0$; then $y > 0$ forces $x > 0$, so $u > 0$ and $v > 0$ throughout $B$.
That is what lets you divide by $x$ and $y$ and rearrange the inequalities without their
directions flipping.
\end{exercise}

\begin{ainote}
The integral in \textbf{8.5.3} is written $\int_C xyz\,dx\,dy$ on the official sheet, but $C$ is
a subset of $\mathbb{R}^3$ and the right-hand side carries three differentials, so this must
read $dx\,dy\,dz$. Quoted verbatim above, as the sheet is a primary source.
\end{ainote}

% Generator: Claude Opus 5 (Medium)
\begin{ainote}
\Cref{ex:8.5} is the direct continuation of \cref{ex:change_of_variables_domain}, and the three
parts are graded. Part 1 is polar coordinates, already computed in
\cref{ex:polar_spherical_coordinates}. Part 3 is a
\emph{linear} change of variables, so $\det\Jac$ is a constant and the transformed domain is
literally the box $(0,1)\times(1,2)\times(-2,0)$ --- read straight off the three constraints.
Part 2 is the interesting one: $1 < xy < 2$ becomes $1 < u < 2$ immediately, but the second
constraint $x^2 < y < 2x^2$ does \emph{not} become a box, since $y = u/\sqrt v$ turns it into
$v^{3/2} < u < 2v^{3/2}$. Straightening one constraint does not oblige the other to cooperate.

Note also that the exercise deliberately stops before the integration: what is being drilled is
finding $\lvert\det\Jac\rvert$ and the new domain, which is where both the marks and the
mistakes are.
\end{ainote}
